% ============================================================
% Capítulo 2 — Arquitectura teórica U-TextReIDNet
% ============================================================
\chapter{U-TextReIDNet: arquitectura teórica}
\label{cap:arquitectura-teorica}

Este capítulo desglosa la arquitectura completa del modelo
U-TextReIDNet tal como se describe en \file{docs/paper.md}
(§\textsc{Iii}-A, líneas 146--178). Aunque el repositorio
\texttt{textreid-train} implementa solo la parte de identificación
(\emph{TextReIDNet}), entender la versión unificada ayuda a
comprender las elecciones de diseño.

\section{Formulación del problema}

Dado un fotograma de vídeo con $n$ personas distintas
$\mathbf{V} \in \mathbb{R}^{H \times W \times 3}$ y una descripción
textual $\mathbf{T}$ en lenguaje natural, los objetivos son:

\begin{enumerate}
  \item \textbf{Detectar y segmentar} las $n$ instancias humanas en
        $\mathbf{V}$, produciendo una galería
        $\mathbf{I} = \{\mathbf{I}_1, \ldots, \mathbf{I}_n\}$.
        Si $n = 0$, la galería está vacía.
  \item \textbf{Identificar} la imagen $\mathbf{I}^* \in \mathbf{I}$
        que mejor se ajusta a $\mathbf{T}$ mediante una función de
        coincidencia multimodal $\mathcal{M}$.
\end{enumerate}

La función $\mathcal{M}$ usa las \emph{incrustaciones} (embeddings)
visuales y textuales:

\begin{align}
  \mathbf{I}' &= F_i(\mathbf{I}) && \text{embedding visual} \\
  \mathbf{T}' &= F_t(\mathbf{T}) && \text{embedding textual}
\end{align}

\noindent y calcula su similitud con la métrica del coseno:

\begin{equation}
  \text{sim}(\mathbf{I}', \mathbf{T}') =
  \frac{\mathbf{I}' \cdot \mathbf{T}'}{\|\mathbf{I}'\|\,\|\mathbf{T}'\|}
  \label{eq:cosine}
\end{equation}

\noindent El resultado óptimo es la imagen que maximiza esa
similitud.

\section{Componentes principales}

U-TextReIDNet se compone de tres bloques:

\begin{enumerate}
  \item \textbf{Red visual}: \emph{backbone} EfficientNet-B0 +
        capa de convolución separable en profundidad.
  \item \textbf{Red de lenguaje}: BERT + BiGRU bidireccional.
  \item \textbf{Componente de coincidencia}: espacio de embedding
        conjunto + pérdidas.
\end{enumerate}

\subsection{Red visual}

\paragraph{Detector R-MHParsNet (no implementado en este proyecto).}

R-MHParsNet es una FPN (Feature Pyramid Network) basada en
EfficientNet-B0, con cabezas de segmentación para generar máscaras
de instancia humana. Genera los bounding boxes
$\{(x_1,y_1,x_2,y_2)\}_{i=1}^{n}$ que delimitan cada persona
detectada. Tiene 6,49M parámetros (un 79,9\% más pequeño que su
predecesor MHParsNet).

\paragraph{Generador de embedding visual (sí implementado).}

Para cada persona detectada, recortada como
$\mathbf{I} \in \mathbb{R}^{384 \times 128 \times 3}$:

\begin{enumerate}
  \item Se redimensiona a $\mathbb{R}^{3 \times 384 \times 128}$
        (HxW o WxH según convención; el paper dice 384 alto × 128
        ancho).
  \item Pasa por EfficientNet-B0 y se extrae la salida de la
        \textbf{etapa 9} (la última): un tensor de
        $\mathbb{R}^{1280 \times 12 \times 4}$. Los 1280 canales
        son la profundidad de features; $12 \times 4$ son las
        dimensiones espaciales (es decir, $384 / 32 = 12$ y
        $128 / 32 = 4$).
  \item Una capa de convolución separable en profundidad
        (\dsc{}) con 1024 filtros $3 \times 3$ reduce a
        $\mathbb{R}^{1024 \times 12 \times 4}$.
\end{enumerate}

\noindent\textbf{Nota técnica:} en este proyecto las imágenes
entran a $512 \times 512$ (ver \fline{datasets/bases.py}{43} y
\cref{cap:dataset}), no a $384 \times 128$. Eso cambia las
dimensiones espaciales en el resto del pipeline. Se conserva la
misma profundidad de canales.

\subsection{Red de lenguaje}

El componente de lenguaje procesa la descripción textual
$\mathbf{T}$:

\begin{enumerate}
  \item \textbf{BERT} tokeniza y recorta/padea a una longitud
        fija de 100 tokens: $(z_{+1}, \ldots, z_{+100})$.
        El tokenizer es \texttt{bert-base-uncased}.
  \item Una \textbf{capa de embedding} de 512 dimensiones convierte
        cada token en un vector denso:
        $(\mathbf{e}_1, \ldots, \mathbf{e}_{100}) \in
        \mathbb{R}^{100 \times 512}$.
  \item Una \textbf{GRU bidireccional} con 1024 unidades ocultas
        procesa la secuencia y produce:
        \begin{equation}
          \mathbf{T}' = (\overrightarrow{\text{GRU}} +
                          \overleftarrow{\text{GRU}}) / 2
          \in \mathbb{R}^{1024 \times 100 \times 1}
          \label{eq:bigru}
        \end{equation}
        La suma promedio de las direcciones hacia adelante y hacia
        atrás captura el contexto en ambos sentidos.
\end{enumerate}

\subsection{Componente de coincidencia de embedding conjunto}

Las dos modalidades se proyectan a un \textbf{espacio conjunto} de
dimensión 1024:

\begin{itemize}
  \item $\mathbf{I}' \in \mathbb{R}^{1024 \times 12 \times 4}$
        (visual).
  \item $\mathbf{T}' \in \mathbb{R}^{1024 \times 100 \times 1}$
        (textual).
\end{itemize}

Para optimizar el \emph{matching} se usan dos pérdidas:

\begin{itemize}
  \item \textbf{Pérdida de identidad} ($\mathcal{L}_{\text{ID}}$):
        cross-entropy que refuerza que cada persona tenga una
        representación única en el espacio.
  \item \textbf{Pérdida de ranking/clasificación}
        ($\mathcal{L}_{\text{rank}}$): basada en triplet loss,
        garantiza que las coincidencias correctas obtengan
        puntuaciones más altas que las incorrectas.
\end{itemize}

\noindent La suma ponderada es la pérdida total:
\begin{equation}
  \mathcal{L} = \alpha \cdot \mathcal{L}_{\text{rank}}
              + \beta  \cdot \mathcal{L}_{\text{ID}}
\end{equation}

\noindent En este proyecto $\alpha = \beta = 1.0$
(\fline{config.py}{87-88}).

\section{Diagrama de bloques}

\begin{figure}[H]
\centering
% ============================================================
% Diagrama TikZ: Arquitectura U-TextReIDNet
% ============================================================
\begin{tikzpicture}[
    font=\sffamily\small,
    node distance=0.6cm and 0.5cm,
    every node/.style={align=center},
    block/.style={
        rectangle, draw=black, thick, rounded corners=2pt,
        minimum height=0.9cm, minimum width=2cm,
        fill=blue!8
    },
    block2/.style={
        rectangle, draw=black, thick, rounded corners=2pt,
        minimum height=0.9cm, minimum width=2cm,
        fill=orange!15
    },
    block3/.style={
        rectangle, draw=black, thick, rounded corners=2pt,
        minimum height=0.9cm, minimum width=2cm,
        fill=green!15
    },
    impl/.style={
        rectangle, draw=blue!60!black, very thick, rounded corners=2pt,
        minimum height=0.9cm, minimum width=2cm,
        fill=blue!15
    },
    arr/.style={
        -{Stealth[length=2.5mm,width=2mm]}, thick
    },
    group/.style={
        rectangle, draw=gray, dashed, thick, rounded corners=4pt,
        inner sep=10pt, fill=gray!5
    }
]

% ======================== RAMA VISUAL ========================
\node[block2] (v_in)  {Frame\\$V \in \mathbb{R}^{H \times W \times 3}$};
\node[block, below=of v_in] (rmhp)  {R-MHParsNet\\(detector)};
\node[block, below=of rmhp] (crop)  {Crop + resize\\$384 \times 128$};
\node[impl, below=of crop] (eff)   {\textbf{EffNet-B0}\\stage 9};
\node[impl, below=of eff]  (dsc1)  {\textbf{DSC}\\1280 $\to$ 1024};
\node[impl, below=of dsc1] (amp)   {Adaptive\\MaxPool 1$\times$1};
\node[impl, below=of amp]  (dsc2)  {\textbf{DSC}\\1024 $\to$ 1024};
\node[impl, below=of dsc2] (ve)    {$\mathbf{I}' \in \mathbb{R}^{1024}$};

% ======================== RAMA TEXTO ========================
\node[block2, right=4cm of v_in] (t_in)  {Texto\\$T$};
\node[impl, right=4cm of rmhp] (bert) {\textbf{BERT}\\100 tokens};
\node[impl, right=4cm of crop] (emb)  {\textbf{Embedding}\\$\mathbb{R}^{512}$};
\node[impl, right=4cm of eff]  (gru)  {\textbf{BiGRU}\\1024 ocultas};
\node[impl, right=4cm of dsc1] (sum)  {Suma fwd+bwd /2};
\node[impl, right=4cm of amp]  (perm) {Permute\\+unsqueeze};
\node[impl, right=4cm of dsc2] (tmax) {Max sobre\\dimensión 2};
\node[impl, right=4cm of ve]   (dsct) {\textbf{DSC}\\1024 $\to$ 1024};
\node[impl, right=4cm of ve, yshift=-0.8cm] (te)  {$\mathbf{T}' \in \mathbb{R}^{1024}$};

% ======================== MATCHING ========================
\node[block3, below=2cm of ve] (cos)  {Similitud coseno\\$\mathcal{M}$};
\node[block3, right=4cm of cos] (rank) {$\mathcal{L}_{\text{rank}}$};
\node[block3, right=4cm of rank] (id)  {$\mathcal{L}_{\text{ID}}$};

% Conexiones rama visual
\draw[arr] (v_in) -- (rmhp);
\draw[arr] (rmhp) -- (crop);
\draw[arr] (crop) -- (eff);
\draw[arr] (eff)  -- (dsc1);
\draw[arr] (dsc1) -- (amp);
\draw[arr] (amp)  -- (dsc2);
\draw[arr] (dsc2) -- (ve);

% Conexiones rama texto
\draw[arr] (t_in) -- (bert);
\draw[arr] (bert) -- (emb);
\draw[arr] (emb)  -- (gru);
\draw[arr] (gru)  -- (sum);
\draw[arr] (sum)  -- (perm);
\draw[arr] (perm) -- (tmax);
\draw[arr] (tmax) -- (dsct);
\draw[arr] (dsct) -- (te);

% Matching
\draw[arr] (ve) -- (cos);
\draw[arr] (te) -- (cos);
\draw[arr] (cos.east) -| (rank.west);
\draw[arr] (cos.east) -| ([yshift=0.6cm]id.west) -- (id.west);
\draw[arr] (rank.east) -- (id.west);

% Agrupaciones
\begin{scope}[on background layer]
  \node[group, fit=(rmhp), label={[gray]above:R-MHParsNet (paper)}] {};
  \node[group, fit=(eff) (dsc1) (amp) (dsc2) (ve), label={[blue!60!black]above:\textbf{Implementado en este proyecto}}] {};
  \node[group, fit=(bert) (emb) (gru) (sum) (perm) (tmax) (dsct) (te), label={[blue!60!black]above:\textbf{Implementado en este proyecto}}] {};
\end{scope}

% Leyenda
\node[draw, rounded corners, fill=orange!15, minimum width=2cm, minimum height=0.5cm, anchor=north west] at (-9, -10) {};
\node[anchor=west] at (-7, -10) {Paper, no en código};
\node[draw, rounded corners, fill=blue!15, minimum width=2cm, minimum height=0.5cm, anchor=north west] at (-9, -10.8) {};
\node[anchor=west] at (-7, -10.8) {Implementado aquí};
\node[draw, rounded corners, fill=green!15, minimum width=2cm, minimum height=0.5cm, anchor=north west] at (-9, -11.6) {};
\node[anchor=west] at (-7, -11.6) {Funciones de pérdida};

\end{tikzpicture}

\caption{Arquitectura de U-TextReIDNet. Las cajas con relleno gris
son las que este proyecto implementa; las cajas con relleno claro
son las del paper que no están en el código (R-MHParsNet y sistema
distribuido).}
\label{fig:arquitectura}
\end{figure}

\section{Dimensiones en cada etapa: tabla resumen}

\begin{table}[H]
\centering
\caption{Dimensiones de los tensores en cada etapa del modelo.
\textit{B} denota el tamaño de batch. Las columnas ``Etapa paper''
y ``Etapa proyecto'' difieren únicamente en la resolución espacial
de entrada.}
\label{tab:dimensiones}
\small
\begin{tabular}{lrll}
\toprule
\textbf{Componente} & \textbf{Forma} & \textbf{Origen} & \textbf{Ubicación} \\
\midrule
Imagen de entrada
  & \dims{(B, 3, 512, 512)}
  & resize/colocar en canvas
  & \fline{datasets/bases.py}{128} \\
\midrule
EfficientNet-B0 stage 9
  & \dims{(B, 1280, 16, 16)}
  & salida \texttt{model.features}
  & \fline{model/visual_network.py}{29} \\
\midrule
DSC 1280 $\to$ 1024
  & \dims{(B, 1024, 16, 16)}
  & conv. separable en profundidad
  & \fline{model/textreidnet.py}{41} \\
\midrule
AdaptiveMaxPool2d(1,1)
  & \dims{(B, 1024, 1, 1)}
  & pool global
  & \fline{model/textreidnet.py}{42} \\
\midrule
DSC 1024 $\to$ 1024
  & \dims{(B, 1024, 1, 1)}
  & conv. separable en profundidad
  & \fline{model/textreidnet.py}{43} \\
\midrule
Visual embedding final
  & \dims{(B, 1024, 1)}
  & \texttt{squeeze(-1)}
  & \fline{model/textreidnet.py}{71} \\
\midrule
\midrule
Tokens de texto (BERT)
  & \dims{(B, 100)}
  & \texttt{bert-base-uncased} + padding
  & \fline{datasets/bert_tokenizer.py}{17} \\
\midrule
Embedding
  & \dims{(B, 100, 512)}
  & \texttt{nn.Embedding(29610, 512)}
  & \fline{model/language_network.py}{22} \\
\midrule
BiGRU
  & \dims{(B, 100, 2048)}
  & bidir, $2 \times 1024$ ocultas
  & \fline{model/language_network.py}{24} \\
\midrule
Suma fwd + bwd / 2
  & \dims{(B, 100, 1024)}
  & reduce a 1024 canales
  & \fline{model/language_network.py}{68} \\
\midrule
Permute + unsqueeze
  & \dims{(B, 1024, 100, 1)}
  & prepara para DSC
  & \fline{model/language_network.py}{69} \\
\midrule
Max sobre dim 2
  & \dims{(B, 1024, 1, 1)}
  & pool temporal
  & \fline{model/textreidnet.py}{87} \\
\midrule
DSC 1024 $\to$ 1024
  & \dims{(B, 1024, 1, 1)}
  & conv. separable en profundidad
  & \fline{model/textreidnet.py}{88} \\
\midrule
Textual embedding final
  & \dims{(B, 1024, 1)}
  & \texttt{squeeze(-1)}
  & \fline{model/textreidnet.py}{88} \\
\bottomrule
\end{tabular}
\end{table}

\section{Por qué estas elecciones}

\subsection*{¿Por qué EfficientNet-B0?}

EfficientNet-B0~\cite{tan2019efficientnet} introduce el concepto de
\emph{compound scaling}: en vez de escalar arbitrariamente
profundidad, ancho o resolución, los escala simultáneamente con un
coeficiente fijo derivado de una búsqueda automática. El resultado
es un modelo con excelente relación accuracy/computo. La variante
B0 es la más pequeña de la familia (\textasciitilde 5,3M parámetros)
y es suficiente para la tarea de extracción de features visuales.

\subsection*{¿Por qué BiGRU en lugar de LSTM o Transformer?}

El paper justifica el uso de \gru{} por su eficiencia computacional:
tiene menos parámetros que LSTM (3 compuertas vs 4) y es más rápido
que un Transformer para secuencias cortas (100 tokens). La
bidireccionalidad le permite capturar el contexto anterior y
posterior, lo que mejora la comprensión de la descripción.

\subsection*{¿Por qué convolución separable en profundidad (DSC)?}

Una convolución separable en profundidad~\cite{howard2017mobilenets}
factoriza una convolución estándar en dos operaciones:

\begin{enumerate}
  \item \textbf{Depthwise}: un filtro $k \times k$ por canal de
        entrada.
  \item \textbf{Pointwise}: una convolución $1 \times 1$ para
        mezclar canales.
\end{enumerate}

\noindent Reduce drásticamente los parámetros y las
multiplicaciones, con una pérdida de accuracy mínima. En este
proyecto se usa \emph{dos} veces: para reducir 1280 $\to$ 1024
canales visuales y para proyectar al espacio conjunto.

\subsection*{¿Por qué 1024 como dimensión de embedding?}

1024 es lo suficientemente grande para capturar la semántica de
las dos modalidades, pero suficientemente pequeño para que la
similitud del coseno sea estable y el entrenamiento sea rápido.
Coincide con el número de unidades ocultas de la BiGRU.

\section{Lo que cambia en este proyecto}

El proyecto \texttt{textreid-train} implementa únicamente el
camino de identificación (\emph{TextReIDNet}) sin el detector
R-MHParsNet. Esto significa:

\begin{itemize}
  \item No se detecta automáticamente a las personas en imágenes
        con múltiples personas: la entrada debe ser una imagen con
        una sola persona (post-recortada o de datasets tipo
        CUHK-PEDES).
  \item Se conserva la \emph{colocación sobre canvas negro}
        (\fline{datasets/bases.py}{43}) porque el paper entrena
        así (las imágenes de CUHK-PEDES son menores que $512
        \times 512$ y se colocan aleatoriamente sobre un lienzo
        negro para preservar aspect ratio).
  \item El tamaño de canvas es $512 \times 512$ en el código
        (no $384 \times 128$ como en el paper). Esto es un punto
        a tener en cuenta si se quiere reproducir exactamente los
        números del paper.
\end{itemize}

Estos detalles se discuten con más profundidad en los capítulos
siguientes.
